% Generated draft: requires mathematical and source review.
\documentclass[11pt]{article}
\usepackage[margin=25mm]{geometry}
\usepackage{fontspec}
\setmainfont{Latin Modern Roman}
\usepackage{amsmath,amssymb,mathtools,unicode-math}
\setmathfont{Latin Modern Math}
\usepackage{enumitem,needspace,xurl,hyperref}
\setlength{\parindent}{0pt}
\setlength{\parskip}{5pt}
\setlength{\emergencystretch}{3em}
\newcommand{\N}{\mathbb{N}}
\newcommand{\Z}{\mathbb{Z}}
\newcommand{\Q}{\mathbb{Q}}
\newcommand{\R}{\mathbb{R}}
\newcommand{\C}{\mathbb{C}}
\newcommand{\F}{\mathbb{F}}
\newcommand{\A}{\mathbb{A}}
\newcommand{\PP}{\mathbb{P}}
\newcommand{\E}{\mathbb{E}}
\newcommand{\eps}{\varepsilon}
\newcommand{\dd}{\,\mathrm{d}}
\begin{document}
\section*{UnsolvedMath ID: 30003611}
\textbf{Source problem code:} catalogue entry
\subsection{Definitions and mathematical statement}
The source leaves the ground field unspecified. The following formulation is conditional on specifying an algebraically closed field \(k\); the intended field and any characteristic restriction remain unresolved.

Let \(X\) be a smooth projective integral variety over \(k\) whose anticanonical divisor \(-K_X\) is ample, so that \(X\) is Fano. Fix an ample Cartier divisor \(L\) on \(X\). Let \(\operatorname{Mor}_k(\PP^1_k,X)\) denote the scheme parametrizing morphisms \(f\colon\PP^1_k\to X\). Define the \(L\)-degree of such a morphism by
\[
\deg_L(f)=\deg_{\PP^1_k}\bigl(f^*\mathcal O_X(L)\bigr).
\]
This degree is constant on each irreducible component \(M\) of the morphism scheme; write it as \(\deg_L(M)\). Interpret a \emph{dominant component} as an irreducible component \(M\) for which the evaluation morphism
\[
\operatorname{ev}_M\colon\PP^1_k\times M\longrightarrow X,
\qquad (t,f)\longmapsto f(t),
\]
has Zariski dense image. This interpretation concerns the family of morphisms parametrized by \(M\), and requires confirmation against the source's intended terminology.

For each integer \(d\geq 0\), set
\[
N_{X,L}(d)=\#\left\{M\subseteq\operatorname{Mor}_k(\PP^1_k,X):\ M\text{ is an irreducible component},\ \operatorname{ev}_M\text{ is dominant},\ \deg_L(M)\leq d\right\}.
\]
The question is whether, for every fixed pair \((X,L)\) over the intended ground field, there exist a real constant \(C>0\) and an integer \(r\geq 0\), independent of \(d\), such that
\[
N_{X,L}(d)\leq C(1+d)^r
\qquad\text{for every integer }d\geq 0.
\]
The constants may depend on \(X\), \(L\), and the ground field; the supplied statement does not request uniformity across varieties or polarizations. The full morphism scheme is counted, including multiple covers and degree-zero morphisms, since the supplied statement makes no exclusion. The Cartier interpretation of the divisor and the meaning of dominant component are explicit conventions for this draft, subject to review.
\subsection{Short English statement}
For a fixed smooth projective Fano variety \(X\) and ample divisor \(L\), count the irreducible components of the scheme of parametrized rational curves whose families cover a dense subset of \(X\) and whose \(L\)-degree is at most \(d\). Must this count grow at most polynomially in \(d\)? The ground field and the precise meaning of dominant component need confirmation. The supplied record labels the problem partially solved; that label is dated curation, not an independently verified conclusion.
\subsection{Sources}
\begin{itemize}
\item Supplied local UnsolvedMath registry record, permanent ID \(30003611\), title ``Polynomial Bounds for Dominant Rational-Curve Components,'' source problem code \texttt{catalogue entry}, in the collection ``Oberwolfach Reports Open Problems.'' The record gives proposed year 2017 and names no individual proposer.
\item Supplied report citation: \emph{Algebraic Geometry: Birational Classification, Derived Categories, and Moduli Spaces} (2018), Oberwolfach Reports, pp.\ 2703--2767, DOI \texttt{10.4171/owr/2017/45}. \url{https://doi.org/10.4171/owr/2017/45}. The supplied background describes these reports as records of mathematical workshops and their open-problem discussions.
\item Supplied literature reference: ``Lehmann Tanimoto 2019.'' \url{https://doi.org/10.1112/S0010437X19007103}.
\item Supplied literature reference: ``Beheshti Lehmann Riedl Tanimoto 2022.'' \url{https://doi.org/10.1016/j.aim.2022.108557}.
\item Literature assessment dated exactly \texttt{2026-08-21}: the supplied curation labels the problem partially solved and states that polynomial component bounds and related uniform estimates are established in broad settings, while the conjectural bound is not proved for all Fano varieties. This is supplied dated literature triage derived from the OWR curation, not independently verified here. The source explicitly requests independent verification of time-sensitive status and cited claims.
\end{itemize}
\end{document}
